% Subsection for Sec. "Quantum analysis" (sec:quantum_analysis) of article_chaos_v2.tex, after the OTOC heatmap.
% New bibliography entries: twa_otoc_refs.bib (to merge into biblio.bib).
% Figure: Chaos/figures_article/otoc_vs_twa_heatmap_22_10_2026.pdf
%   produced by plot_OTOCS.ipynb (cell baed0266) from twa_otoc_heatmap_22_10_2026.json (twa_otoc_heatmap.py).
% Late-time data: otoc_attractor_start_na65_nb33_22_10_2026.json, otoc_stochastic_twa_22_10_2026.json.

\subsection{Classical counterpart of the OTOC: truncated Wigner approximation}
\label{subsec:twa}

The onset of positive OTOC rates in \cref{fig:OTOC_coherent} does not follow the classical onset of chaos. It sits at $g_2/2\pi\simeq100~\mathrm{MHz}$ for all cat sizes, whereas the onset of the Lyapunov map started on a cat lobe moves from $117$ to $78~\mathrm{MHz}$ as $|\alpha|^2$ grows from $3$ to $7$ (\cref{fig:lyapunov_heatmap_cat_start}). Two effects could explain the difference: quantum effects absent from the classical dynamics, or a difference between the two observables. The OTOC is a finite-time quantity, averaged over a wave packet that starts away from the attractor, while the Lyapunov exponent is a long-time property of a single trajectory on it. To separate the two effects, we compute the classical quantity that corresponds term by term to the OTOC, with the truncated Wigner approximation (TWA)~\cite{Steel1998Dec,Sinatra2002Jun,Polkovnikov2010Sep}.

\paragraph*{Truncated Wigner approximation.}
In the Wigner representation $W(\alpha,\beta,t)$ of the two modes, the GKSL equation \cref{eq_app_GKSL} becomes a partial differential equation. Its first-order derivative terms are the Liouville flow of the classical equations \cref{eq_classical_dynamics}, with $\cos X$ evaluated at each point of phase space. The single-photon loss adds the drift $-(\kappa_b/2)\beta$ and the diffusion $(\kappa_b/2)\,\partial_\beta\partial_{\beta^*}W$, the vacuum noise of the loss. The nonlinearity also generates third- and higher-order derivatives, whose weight relative to the flow is set by the zero-point fluctuations of the junction phase, $\varphi_a^2$ and $\varphi_b^2\lesssim0.04$. The TWA drops them. What remains is a Fokker--Planck equation, equivalent to the stochastic equations
\begin{equation}
\begin{split}
    \dot b &= F_b(a,b,t)+\sqrt{\kappa_b/2}\;\xi(t),\\
    \langle\xi(t)\xi^*(t')\rangle &= \delta(t-t'),
\end{split}
\label{eq_twa_langevin}
\end{equation}
where $F_b$ is the right-hand side of the equation for $\dot b$ in \cref{eq_classical_dynamics}, with $\dot a$ unchanged and the initial conditions sampled from the Wigner function of the initial state. For a coherent state $\ket{a_0,b_0}$, the samples are Gaussian around $(a_0,b_0)$ with $\langle|\delta a|^2\rangle=\langle|\delta b|^2\rangle=1/2$. The noise relaxes the bare damped buffer to the vacuum, $\langle|\delta b|^2\rangle=1/2$. The TWA is exact at first order in the quantum fluctuations and holds up to the Ehrenfest time, at which the neglected terms become comparable to the flow; in a chaotic system this time grows only logarithmically with the inverse effective Planck constant~\cite{Polkovnikov2010Sep}. As for the Lyapunov maps, we integrate the equations in the quadratures of \cref{eq:app_lyap_eom}, with the noise as a force $\sqrt{\kappa_b}\,\omega_b\,\xi(t)$, $\langle\xi(t)\xi(t')\rangle=\delta(t-t')$, added to $\ddot x_b$. For $\kappa_b\ll\omega_b$ it is equivalent to the noise of \cref{eq_twa_langevin}, as the damping $-\kappa_b\dot x_b$ is to $-(\kappa_b/2)b$, and it relaxes the bare buffer to the vacuum, $\langle x_b^2\rangle=1/2$ ($0.49$ numerically).

\paragraph*{OTOC in the truncated Wigner approximation.}
At leading order in the quantum fluctuations, the commutator becomes a Poisson bracket, $-i[\hat x_a(t),\hat p_a(0)]\to\{x_a(t),p_a(0)\}=\partial x_a(t)/\partial x_a(0)$, and the expectation value becomes an average over the Wigner function. The OTOC thus becomes
\begin{equation}
    C_\mathrm{TWA}(t)=\Bigl\langle\Bigl(\frac{\partial x_a(t)}{\partial x_a(0)}\Bigr)^2\Bigr\rangle ,
    \label{eq_otoc_twa}
\end{equation}
where $x_a(t)$ is the position in the frame rotating with the memory, as in the quantum calculation, and the average runs over the initial Wigner samples and the noise~\cite{Schmitt2019Apr,Rammensee2018Sep}. The derivative follows the tangent dynamics along each stochastic trajectory. Since $\partial x_a(t)/\partial x_a(0)\sim e^{\Lambda_t t}$, with $\Lambda_t$ the finite-time Lyapunov exponent of that trajectory, \cref{eq_otoc_twa} is an average of exponentials, $C_\mathrm{TWA}=\langle e^{2\Lambda_t t}\rangle\ge e^{2\langle\Lambda_t\rangle t}$. Its growth rate therefore exceeds the Lyapunov exponent and is set by the most unstable trajectories of the initial cloud~\cite{Rozenbaum2017Feb}. It can even be positive near unstable fixed points of a regular dynamics~\cite{Xu2020Apr}. Moreover, over $\kappa_bT=0.225$ from $\ket{i\alpha,0}$, which is not on the attractor (the buffer alone takes $1/\kappa_b$ to fill), the rate measures the transient stretching of the initial cloud and not the asymptotic exponent. The Lyapunov exponent of \cref{fig:lyapunov_heatmap_cat_start} is recovered from $\partial x_a(t)/\partial x_a(0)$ only along a single trajectory, after the transient and over times of tens of nanoseconds. We checked this with the same tangent dynamics, started on the cat lobe.

\paragraph*{Comparison.}
We evaluate \cref{eq_otoc_twa} on the grid, initial state and time window of the quantum OTOC of \cref{fig:OTOC_coherent} and fit it in the same way. \Cref{fig:otoc_vs_twa} shows the two maps side by side. They agree in sign on $288$ of the $299$ points. Their zero lines coincide within one grid step (about $3~\mathrm{MHz}$) for every cat size except the smallest, $|\alpha|^2=1$, where they differ by two. Both rates are negative and almost equal in the regular region, and they reach comparable values deep in the chaotic region. The mean of exponentials is dominated by rare trajectories: we used $16384$ Wigner samples, which converges the rates to about $0.3\,\kappa_b$. On this window, the quantum OTOC is therefore the classical dynamics of the initial Wigner distribution. Its difference with the Lyapunov maps comes from the observable, not from quantum effects.

\begin{figure*}[tbp]
    \centering
    \includegraphics[width=\textwidth]{Chaos/figures_article/otoc_vs_twa_heatmap_22_10_2026.pdf}
    \caption{Growth rate of the OTOC from the coherent state $\ket{i\alpha,0}$, $|\alpha|^2=\eps_d/2g_2$, over three drive periods ($\kappa_bT=0.225$): half the slope of a linear fit of $\ln C$ against $\kappa_bt$. (a)~Quantum OTOC, \cref{app:otoc,app:OTOC_coherent}. (b)~Truncated Wigner approximation, \cref{eq_otoc_twa}, with $16384$ samples of the Wigner function of the initial state and the vacuum noise of the buffer loss. The solid and dashed lines are the zero lines of the quantum and TWA rates. Same color scale in both panels.}
    \label{fig:otoc_vs_twa}
\end{figure*}

\paragraph*{Beyond the Ehrenfest time.}
To probe the chaotic attractor itself, we also started the quantum OTOC from coherent states centred on points of classical trajectories taken after the transient of the Lyapunov maps, at $\eps_d=10g_2$ and $g_2/2\pi$ from $78$ to $110~\mathrm{MHz}$, over $24$ drive periods ($\kappa_bT=1.8$). The OTOC follows $C_\mathrm{TWA}$ for the first one to four drive periods. It then saturates at $C\simeq15$--$40$ in the chaotic region (about $5$ in the regular region) after $0.6$ to $1~\mathrm{ns}$, and decays back to $C\simeq1$ by $2$ to $3~\mathrm{ns}$, while $C_\mathrm{TWA}$ keeps growing. The saturation at the Ehrenfest time is generic~\cite{Rozenbaum2017Feb,Rammensee2018Sep}. The decay is specific to the open system. The Heisenberg operator $\hat x_a(t)=\Phi^\dagger(t,0)\hat x_a$ relaxes towards $\langle\hat x_a\rangle_\mathrm{ss}\mathbb 1$ when the steady state is unique, so the commutator vanishes~\cite{Zanardi2021Jun,Touil2021Jan}. Classically, this corresponds to averaging the derivative over the noise before squaring, $\langle(\partial_{x_a(0)}\mathbb E_\xi[x_a(t)])^2\rangle$, since $\Phi^\dagger$ traces out the bath. This ordering reproduces the saturation and the decay at the same times, with a peak two to three times lower, and without the noise the decay disappears. It underestimates the growth at short times, however, which the mean of squares \cref{eq_otoc_twa} captures. The quantum OTOC thus crosses over from one ordering to the other around the Ehrenfest time, a regime beyond the leading order of the TWA. A Lyapunov exponent can therefore not be extracted from the OTOC of this open system: its exponential growth lasts less than an Ehrenfest time, about one nanosecond here.
\TMD{The crossover between the two orderings is our interpretation; the data are in otoc\_attractor\_start\_na65\_nb33\_22\_10\_2026.json and otoc\_stochastic\_twa\_22\_10\_2026.json (plot\_OTOCS.ipynb, cell e7634ee0). A figure could go to an appendix.}
